\section{The Represented Stokes Theorem}
\label{sec:represented-stokes}

The central statement of this artifact is a represented Stokes theorem.  The
word represented is meant in the same mathematical spirit as a represented
measure, a coordinate expression, or a chain-level representative: the theorem
chooses a canonical finite coordinate expression produced by the proof and
proves Stokes at that level.  It is not a statement about approximation, and it
is not a placeholder for a missing equality.

\subsection{The classical finite-coordinate endpoint}

Let \(\omega\) be a compactly supported \(n\)-form on a manifold modeled on a
half-space.  The usual proof of Stokes begins by choosing a finite family of
charts and a partition of unity \((\rho_j)\) on the compact support.  The bulk
side is decomposed formally as
\[
  \int_M \dd\omega
  =
  \sum_j \int_M \dd(\rho_j\omega),
\]
up to the standard partition and support arguments.  For interior charts, the
coordinate theorem reduces the local term to an ordinary box theorem whose
boundary contributions are artificial and vanish by support.  For boundary
charts, each localized term is further refined into finitely many half-space
boxes.  On such a box \(Q_{i,q}\), the local theorem gives a term of the form
\[
  B_{i,q}(\omega) = C_{i,q}(\omega),
\]
where \(B_{i,q}\) is the local bulk integral of the coordinate representative
of a localized form and \(C_{i,q}\) is the true lower-face boundary contribution
with the half-space boundary sign.

The represented bulk endpoint is the nested finite sum of the \(B_{i,q}\).  The
represented boundary endpoint is the nested finite sum of the \(C_{i,q}\).  In
schematic form,
\[
  \BulkRep(\omega)
    = \sum_i \sum_{q \in Q_i} B_{i,q}(\omega),
  \qquad
  \BoundaryRep(\omega)
    = \sum_i \sum_{q \in Q_i} C_{i,q}(\omega).
\]
The theorem proves \(\BulkRep(\omega)=\BoundaryRep(\omega)\) for the canonical
finite data generated from compact support and chartwise smoothness.

\begin{definition}[Represented Stokes route]
A represented Stokes route for a compactly supported form consists of finite
boundary chart indices \(i\), finite refined half-space pieces \(q\in Q_i\),
box corners \(a_{i,q}\le b_{i,q}\) with \((a_{i,q})_0=0\), localized
coefficients \(\rho_{i,q}\), and coordinate representatives
\(\alpha_{i,q}\) satisfying the following conditions.
\begin{enumerate}
  \item The coefficients arise from a support-controlled partition and smooth
  refinement of the compact carrier.
  \item The topological support of \(\alpha_{i,q}\) is contained in the
  assigned half-space support box, so every artificial face of the coordinate
  box gives zero.
  \item The representative \(\alpha_{i,q}\) satisfies the open-interior and
  closed-box regularity assumptions needed for the half-space local theorem.
  \item The local bulk and lower-face boundary terms assemble into generated
  finite endpoint sums.
\end{enumerate}
The represented endpoints of the route are the two finite sums produced by the
last item.
\end{definition}

The Lean construction makes this definition concrete.  The endpoint layer
stores the refined pieces and local terms in
\code{CoverIndexedBoundaryLocalizedRefinedPartition}; the zero-localized
certificate supplies the support condition; and the selected regularity
constructor supplies the analytic fields.  The first-principles theorem proves
that the route exists canonically from the theorem hypotheses.

\subsection{Endpoint algebra}

The represented endpoint is best understood as a finite algebraic object.  For
each active pair \((i,q)\), the route assigns a localized coefficient
\(\rho_{i,q}\), a boundary chart \(x_i\), and half-space box corners
\((a_{i,q},b_{i,q})\).  The local bulk term has the form
\[
  B_{i,q}
  =
  \int_{[a_{i,q},b_{i,q}]}
    \operatorname{coeff}_{\mathrm{top}}
      \bigl(\dd(\operatorname{ZeroRep}_{x_i,x_i}(\rho_{i,q}\omega))\bigr),
\]
where the integral notation abbreviates the coordinate box integral used by
the formal local theorem.  The boundary term is
\[
  C_{i,q}
  =
  \operatorname{sgn}_{\partial\HH}
  \int_{\{x_0=0\}\cap [a_{i,q},b_{i,q}]}
    \operatorname{coeff}_{\partial}
      \bigl(\operatorname{ZeroRep}_{x_i,x_i}(\rho_{i,q}\omega)\bigr).
\]
The represented endpoints are then exactly
\[
  \BulkRep(\omega)=\sum_i\sum_{q\in Q_i} B_{i,q},
  \qquad
  \BoundaryRep(\omega)=\sum_i\sum_{q\in Q_i} C_{i,q}.
\]

This display is schematic, but it captures the formal structure: both
endpoints use the same finite selected pieces, the same refined coefficients,
and the same chart choices.  The local Stokes theorem proves \(B_{i,q}=C_{i,q}\)
piecewise; endpoint reconstruction proves that summing these equalities gives
the equality of the generated endpoint definitions.

This is exactly the equality obtained in the middle of the textbook proof,
after localization and before comparison with a global manifold integration
interface.  In informal mathematics one usually identifies this finite
coordinate expression with the left and right sides of the global theorem
immediately.  In Lean, that comparison should be a theorem about the chosen
global integration API.  The present work isolates and proves the coordinate
Stokes content independently of that future comparison layer.

\subsection{What is generated, and what is not supplied}

The strongest evidence that the represented endpoints are mathematical rather
than arbitrary is their construction.  In the final route, the user does not
provide the endpoints.  The top-level structure
\[
\code{CoverIndexedZeroCompactRepresentedStokesFirstPrinciplesInput}
\]
has two fields:
\begin{lstlisting}[language=Lean]
structure CoverIndexedZeroCompactRepresentedStokesFirstPrinciplesInput
    (I : ModelWithCorners Real (Fin (n + 1) -> Real) H)
    (omega : ManifoldForm I M n) where
  chartwiseSmooth : ManifoldForm.ChartwiseSmooth I omega
  compactSupport : IsCompact (closure (ManifoldForm.support I omega))
\end{lstlisting}
The canonical bulk and boundary endpoints are definitions in the namespace of
this structure:
\begin{lstlisting}[language=Lean]
def canonicalRepresentedBulkIntegral : Real := ...
def canonicalRepresentedBoundaryIntegral : Real := ...
\end{lstlisting}
They are obtained by first building an intrinsic input on the compact carrier
\(\overline{\supp(\omega)}\), then following the canonical selected-cover and
refinement route.  The theorem is:
\begin{lstlisting}[language=Lean]
theorem representedStokes :
    X.canonicalRepresentedBulkIntegral (I := I) (omega := omega) =
      X.canonicalRepresentedBoundaryIntegral (I := I) (omega := omega)
\end{lstlisting}

This shape matters.  A certificate-consuming theorem would have asked the user
to provide a finite cover, local Stokes data, endpoint adapters, and finally
the names of the two real numbers to compare.  The first-principles theorem
does the opposite.  It internalizes the choices and exposes the generated
endpoint names as projections of the proof route.

The noncomputable choices inside the proof do not weaken the statement.  They
are the ordinary choices already present in the classical proof: finite
subcovers and subordinate partitions.  The theorem fixes a construction path
and proves the equality for the endpoints generated by that path.  A future
native-integral comparison theorem can then prove that these generated
endpoints agree with a chart-independent integral, but the Stokes equality of
the generated endpoints is already a theorem.

\subsection{The represented-to-native boundary}

The represented theorem should be read as a theorem at a stable mathematical
boundary.  It proves the finite coordinate equality produced by the classical
proof.  It does not yet assert:
\[
  \BulkRep(\omega) = \int_M \dd\omega,
  \qquad
  \BoundaryRep(\omega) = \int_{\partial M} \omega
\]
for a particular future or external native manifold integration API.  Such a
comparison requires a convention for integrating differential forms over
oriented manifolds with boundary in mathlib, boundary orientation data, measure
normalization, chart-independence theorems, and partition-independence
theorems.

The distinction is useful rather than apologetic.  It separates two logically
different parts of Stokes:
\begin{itemize}
  \item the local-to-global Stokes proof, where compact support, charts,
  partitions, half-space boxes, artificial faces, and signs are handled;
  \item the representation theorem identifying the resulting finite coordinate
  sums with a native manifold integral interface.
\end{itemize}
The artifact formalizes the first part at scale.  The second part is the
natural next interface once a suitable native integration API is available.

\subsection{Why represented endpoints are canonical}

There are several sources of noncomputability in the theorem: finite
subcovers, partitions of unity, and choices of refinements.  This does not make
the endpoint arbitrary.  In Lean, ``canonical'' here means canonical relative
to the chosen construction path in the theorem namespace.  The endpoint is
defined by that path, and the theorem proves equality for those generated
definitions.  This is analogous to choosing a proof-generated representative
in homological algebra and proving a chain-level equality before proving
independence of representative.

The formal route fixes the carrier \(K\), extracts a selected open cover from
pointwise chart-box data, constructs a support-controlled partition, refines
boundary active carriers, builds local zero-localized Stokes data, and then
defines the endpoints by finite sums over the resulting boundary pieces.  Thus
the represented values are stable enough to state, audit, and use as the target
of future native comparison theorems.

\begin{theorem}[Represented local-to-global Stokes, mathematical reading]
Let \(\omega\) be a chartwise smooth compactly supported \(n\)-form in the
half-space model context of the artifact.  For the finite represented Stokes
route generated from \(\overline{\supp(\omega)}\), every refined local
half-space term satisfies local Stokes, artificial coordinate faces vanish, and
therefore the generated represented endpoints satisfy
\[
  \BulkRep(\omega)=\BoundaryRep(\omega).
\]
\end{theorem}

\begin{definition}[Coordinate-represented compact-support Stokes]
In this paper, a coordinate-represented compact-support Stokes theorem is a
Lean theorem which, from chartwise smoothness and compactness of the closed
support, constructs finite coordinate data following the standard local proof
and proves equality of the generated finite bulk and boundary coordinate
endpoints.
\end{definition}

This definition is the claim made by the artifact.  It is deliberately stronger
than a local Euclidean theorem and deliberately narrower than a full
mathlib-native manifold integral statement.
